\section{Limitations}
\label{sec:limitations}

% NOTE: polished your four bullet points into prose; item 1's original wording
% ("which our tokenizer can not do shit") is toned down below to "the
% tokenizer cannot compensate for this" -- flagging the rewrite explicitly in
% case I changed the intended meaning, not just the register.

The method inherits several limitations directly from its design.

\textbf{Graph structure constraints.} The context tree and the coverage
score (Section~\ref{sec:method}) can only ever reflect the relationships that
are actually encoded in the graph. Where SNOMED~CT's own definition of a
concept is incomplete -- a relevant attribute simply not modeled as an
outgoing edge -- the tokenizer cannot compensate for this: an under-specified
concept will look under-specified in its context tree regardless of which
token set $T$ is chosen. \doubt{"which our tokenizer can not do shit" in your
original note -- I read this as "our method has no way to fix or detect
this," which is what the sentence above says; flag if you meant something
more specific, e.g.\ a particular class of missing relationships you've
already run into.}

\textbf{Transductive selection.} Token selection (Section~\ref{sec:method:selection})
is transductive: it optimizes coverage for a fixed, known $M$, and a token
set selected for one institution's mapped concepts is not guaranteed to
transfer to a different mapped set without re-running selection.

\textbf{No external validation on a clinical task.} \doubt{This point is in
tension with two other places in your outline: Method item 5 lists a
"(maybe)" clinical-expert acceptability evaluation, and Future Work item 2
lists "external validation clinical task" as planned. As drafted here I've
kept this limitation as stated (no external validation exists in this
paper), but if the phenotype-discovery / IA-cohort case study
(\texttt{phenotype\_clustering.py}, still live in \texttt{src/} even though
its driving notebooks sit in the now-obsolete \texttt{old\_notebooks/}) is
meant to be part of this paper, this limitation needs to be softened or
removed rather than stated as-is -- please confirm whether that case study is
in scope.}

\textbf{Requires a DAG.} The tokenizer's core mechanism -- \texttt{IS\_A}
transparency, non-\texttt{IS\_A} subcontexts, and the coverage recursion's
reliance on well-defined finite hop distances -- assumes an acyclic graph. It
does not extend as-is to a vocabulary with cycles, or to a general knowledge
graph where the child/parent asymmetry of Section~\ref{sec:background} does
not hold.
